The autonomous capture of non-cooperative targets by Unmanned Aerial Vehicles (UAVs) faces critical challenges due to the semantic gap of classical vision algorithms and the computational latency that destabilizes the aircraft. To overcome these limitations, this work develops a real-time Position-Based Visual Servoing (PBVS) system for the approach and capture of objects using an underside gripper. The architecture integrates deep learning models from the YOLO family for robust, marker-less target detection. To mitigate occlusions and noise caused by the rotors, an Extended Kalman Filter (EKF) acts in the image plane as a kinematic predictor. The two-dimensional coordinates are converted to metric three-dimensional space via the Pinhole model, feeding a cascaded Proportional-Derivative (PD) flight controller responsible for compensating visual latency. The system was modeled in a Software-in-the-Loop (SITL) simulation environment using ROS and Gazebo. The project targets a visual processing rate above 20 FPS, a 20\% reduction in trajectory error (RMSE), overshoot maintenance below 15\%, and a final approach precision of 0.30 meters, validating a viable and stable solution for autonomous aerial manipulation.